\documentclass[a4paper,12pt]{report}
\usepackage{longtable}
\usepackage[left= 2.5 cm,right = 2.5 cm, bottom = 2 cm]{geometry}
\usepackage[utf8]{inputenc}
\usepackage[T1]{fontenc}
\usepackage[ngerman]{babel}
\usepackage{titlesec, graphicx, hyperref, chngcntr, url, changepage, ragged2e, setspace, mathptmx, courier, caption}
\usepackage[printonlyused]{acronym}
\usepackage[svgnames]{xcolor}
\usepackage[parskip]{}
\usepackage[scaled=0.9]{helvet}
\usepackage{amsmath}
\usepackage{fancyhdr}
\usepackage{abstract}
\usepackage{lmodern}
\usepackage{adjustbox}
%Graphen Plotten
\usepackage{pgfplots}
\usepackage[style=ieee,
backend=biber, natbib=true]{biblatex}
\addbibresource{Literatur.bib}
\pgfplotsset{compat = newest}

\pagestyle{fancy}
\fancyhf{}
\fancypagestyle{plain}{%  the preset of fancyhdr 
    \fancyhf{} % clear all header and footer fields
    \fancyfoot[C]{\textbf{\thepage}} % except the center
    \fancyhead[C]{\leftmark}
    \renewcommand{\headrulewidth}{0.6pt}
    \renewcommand{\footrulewidth}{0.4pt}}

% Generelle Einstellungen
% Hurenkinder und Schusterjungen verhindern
\clubpenalty = 10000 % schließt Schusterjungen aus (Seitenumbruch nach der ersten Zeile eines neuen Absatzes)
\widowpenalty = 10000 % schließt Hurenkinder aus (die letzte Zeile eines Absatzes steht auf einer neuen Seite)
\displaywidowpenalty=10000
\urlstyle{same}
\definecolor{accent_color}{HTML}{104EBC}
\titleformat{\chapter}{\LARGE\bfseries}{\thechapter\hspace{10pt}\textcolor{accent_color}{|}\hspace{10pt}}{0pt}{\LARGE\bfseries}
\titlespacing*{\chapter}{0pt}{-15pt}{20pt}
\setlength{\parskip}{1em}

% Farbe der Zitate anpassen
\hypersetup{hidelinks}

\begin{document}

    % Deckblatt
    \begin{spacing}{1}
        %!TEX root = ../Studienarbeit.tex
\begin{titlepage}
	\begin{longtable}{p{8.2cm} p{5.4cm}}
        {\raisebox{\ht\strutbox}{\includegraphics[height=2.5cm]{Bilder/TruLogo.jpg}}} &
		{\raisebox{\ht\strutbox}{\includegraphics[height=2.5cm]{Bilder/dhbw.png}}}
	\end{longtable}
	\enlargethispage{20mm}
	\begin{center}
		\vspace*{12mm}	{\LARGE\textbf{{Datengenerator für Daten mit Bias als}}}\\
		\vspace*{3mm} {\LARGE\textbf{{Grundlage für Data Science Projekte}}}\\
		\vspace*{12mm}	{\large\textbf{Studienarbeit}}\\
		\vspace*{12mm}	für die Prüfung zum\\
		\vspace*{3mm}		{\textbf{Bachelor of Science}}\\
		\vspace*{12mm}	des Studiengangs Informatik\\
    \vspace*{3mm}		an der Dualen Hochschule Baden-Württemberg Stuttgart\\
		\vspace*{12mm}	von\\
		\vspace*{3mm}		{\large\textbf{Simon Jess, Timo Zaoral}}\\
		\vspace*{12mm}	Juni 2022\\
	\end{center}
	\vfill
	\begin{spacing}{1.2}
	\begin{tabbing}
		mmmmmmmmmmmmmmmmmmmmmmmmmm             \= \kill
		\textbf{Bearbeitungszeitraum}       \>  04.10.2021 - 10.06.2022\\
		\textbf{Matrikelnummer, Kurs}  \>  8268544, 6146532, INF19C\\
		\textbf{Ausbildungsfirma}                  \>  TRUMPF GmbH \& Co KG, Ditzingen\\
		\textbf{Betreuer}               \>  Prof. Dr. Monika Kochanowski \\ \\ \\
	\end{tabbing}
	\end{spacing}
\end{titlepage}
        \phantomsection
    \end{spacing}
    \newpage

	% Sperrvermerk
	%\input{ads/sperrvermerk}
	%\newpage

	% Erklärung
	\input{ads/erklärung}
    \phantomsection
	\newpage

    %%% ABSTRACT
    \renewcommand{\abstractname}{Abstract}
    \begin{abstract}
        Fasst die Aufgabenstellung und Ergebnisse kompakt und übersichtlich in wenigen Zeilen zusammen (4-7 Zeilen).
        \thispagestyle{empty}
    \end{abstract}
    \phantomsection
    %%% ENDE ABSTRACT

	% Inhaltsverzeichnis
	\begin{spacing}{1}
		\begingroup
            \fancyhf{}
            \fancypagestyle{plain}{%  the preset of fancyhdr 
                \fancyhf{} % clear all header and footer fields    
                \renewcommand{\footrulewidth}{0pt}}
                \renewcommand{\headrulewidth}{0pt}
			\setcounter{tocdepth}{3}
			\tableofcontents
			\clearpage
		\endgroup
	\end{spacing}
    \phantomsection

    \pagenumbering{Roman}
    \setcounter{page}{5}
    \fancyhf{}
	\fancypagestyle{plain}{%  the preset of fancyhdr 
        \fancyhf{} % clear all header and footer fields
        \fancyfoot[C]{\textbf{\thepage}} % except the center      
        \renewcommand{\footrulewidth}{0.4pt}}
        \renewcommand{\headrulewidth}{0pt}
    \newpage

    %Abkürzungsverzeichnis
    \thispagestyle{empty}
    \input{ads/abkürzungen}
    \newpage

    % Abbildungsverzeichnis
	\cleardoublepage
    \thispagestyle{empty}
    \phantomsection
	\listoffigures
    \addcontentsline{toc}{section}{Abbildungsverzeichnis}
    \newpage

	%Tabellenverzeichnis
	\cleardoublepage
    \thispagestyle{empty}
    \phantomsection
	\listoftables
    \addcontentsline{toc}{section}{Tabellenverzeichnis}
    \newpage

    \pagenumbering{arabic}
	\fancyhf{}
    \fancypagestyle{plain}{%  the preset of fancyhdr 
        \fancyhf{} % clear all header and footer fields
        \fancyfoot[C]{\textbf{\thepage}} % except the center
        \fancyhead[C]{\leftmark}
        \renewcommand{\headrulewidth}{0.6pt}
        \renewcommand{\footrulewidth}{0.4pt}}
    \pagestyle{plain}
    % Inhalt
	\foreach \i in {01,02,03,04,05,06,07,08,09,...,99} {%
    \edef\FileName{content/\i kapitel}%
        \IfFileExists{\FileName}{%
            \input{\FileName}
        }
        {%
            %file does not exist
        }
    }
    
    %Literaturverzeichnis
    \newpage
    \printbibliography
    %\bibliographystyle{ieeetr}
    %\bibliography{Literatur.bib}

\end{document}